\section{Shell paths and gradings}\label{sec:shells}
\subsection{Shell constructions and periodic growth}
\begin{definition}\label{def:shell}
Let $j\in\Z_{\ge0}$ and $\epsilon\in\Z_{\ge0}$ satisfy $j+1+\epsilon\le m$.
For $s\in[j]_{\Z}$, let $X_s$ be a tNE path.
Let $u\in[m-1]_{\Z}$ and $v\in[m-1]_{\Z}$.
Let $X=X_1\ccn \cdots\ccn X_j$.
For $j=0$, let $X=\emptyset$.
\begin{enumerate}[font=\itshape]
\item We define the \textit{shell path}
\[
 Y=\operatorname{Shell}(X_1,\ldots,X_j;u,v,\epsilon)=R_u\ccn X\ccn R_v.
\]
For $j=0$, use the notation $\operatorname{Shell}(;u,v,\epsilon)$.
\item For $s\in[j]_{\Z}$, let $g_s\in\operatorname{CtdG}_m(X_s)$.
Let $f$ be the concatenated grading on $X$.
For $j=0$, let $f=\emptyset$.
We define the \textit{shell grading} $g=\operatorname{GShell}_{f}(X_1,\ldots,X_j;u,v,\epsilon)$ on $Y$ by
\begin{equation}\label{eq:shell}
 g(Y_k)=
 \begin{cases}
 u+1,&k=1,\\
 f(X_{k-u-1}),&k\in[u+2,u+1+|X|]_{\Z},\\
 v-1,&k=u+|X|+2,\\
 j+1+\epsilon,&k=|Y|,\\
 0,&\text{otherwise}
 \end{cases}
 \qquad(k\in[|Y|]_{\Z}).
\end{equation}
For $j=0$, use the notation $\operatorname{GShell}_{\emptyset}(;u,v,\epsilon)$.
\item We call a shell path $\operatorname{Shell}(X_1,\ldots,X_j;u,v,\epsilon)$ \textit{unary} if $j=1$.
\item We call a shell path $\operatorname{Shell}(X_1,\ldots,X_j;u,v,\epsilon)$ \textit{binary} if $j=2$.
\item We call a shell grading $\operatorname{GShell}_{f}(X_1,\ldots,X_j;u,v,\epsilon)$ \textit{preserving} if $\epsilon=0$.
\item We call a shell grading $\operatorname{GShell}_{f}(X_1,\ldots,X_j;u,v,\epsilon)$ \textit{draining} if $j=\epsilon=1$.
\end{enumerate}
\end{definition}

\begin{definition}\label{def:recolor}
Let $U$ and $V$ be tNE paths. Let $S\subseteq\operatorname{NE}(U)$ and
$\iota\colon S\hookrightarrow\operatorname{NE}(V)$.
Let $g\in G_m(U)$, $\gamma\in\operatorname{Col}(U,g,m)$, and $h\in G_m(V)$.
We define the \textit{recoloring operator}
\[
 \operatorname{Recol}(U,V,\iota,h;g,\gamma)=\gamma'\in\operatorname{Col}(V,h,m)
\]
such that $\gamma'$ satisfies
\begin{equation}\label{eq:color-update}
 \gamma'(e)=
 \begin{cases}
 \gamma(x),&e=\iota(x),\ x\in S,\ h(e)=g(x),\\
 1,&\text{otherwise}
 \end{cases}
 \qquad(e\in\operatorname{NE}(V)).
\end{equation}
\end{definition}
\Needspace{7\baselineskip}
\begin{lemma}\label{lem:shell}
Let $j\in\Z_{\ge0}$ and $\epsilon\in\Z_{\ge0}$ satisfy $j+1+\epsilon\le m$. Let $u\in[m-1]_{\Z}$ and $v\in[m-1]_{\Z}$.
For $s\in[j]_{\Z}$, let $X_s$ be a tNE path and let $g_s\in\operatorname{CtdG}_m(X_s)$.
Let $X=X_1\ccn\cdots\ccn X_j$ and let $f$ be the concatenated grading on $X$.
For $j=0$, let $X=\emptyset$ and $f=\emptyset$.
Suppose that one of the following conditions holds.
\begin{hypotheses}
\item\label{hyp:shell-preserving} $j\in[0,2]_{\Z}$ and $\epsilon=0$.
\item\label{hyp:shell-draining} $j=\epsilon=1$ and $\operatorname{Surp}(X_1,g_1)\ge1$.
\end{hypotheses}
Let $Y=\operatorname{Shell}(X_1,\ldots,X_j;u,v,\epsilon)$ and $g=\operatorname{GShell}_{f}(X_1,\ldots,X_j;u,v,\epsilon)$.
Put
\[
 r_i=u+1+\sum_{s=1}^i|X_s|\qquad(i\in[0,j]_{\Z}).
\]
The empty sum is $0$.
Then, we have the following.

\begin{enumerate}
\item\label{c:shell-bounded} $g\in G_m(Y)$.
\item\label{c:shell-connected} $g\in\operatorname{CtdG}_m(Y)$.
\item\label{c:shell-residual} We have
\[
 \operatorname{Surp}(Y,g)=\sum_{i=1}^j\operatorname{Surp}(X_i,g_i)-\epsilon.
\]
\item\label{c:shell-recolor} Fix $X_1,\ldots,X_j$, $u$, $v$, and $\epsilon$.
Let $c_s\in\Z_{\ge0}$ for $s\in[j]_{\Z}$.
Assume $\sum_{s=1}^j c_s\ge\epsilon$.
For $(g_s,\gamma_s)\in\mathcal C_{c_s}(X_s)$, let $\eta$ be the concatenated coloring.
If $j=0$, let $\eta=\emptyset$.
We define
\[
 \iota\colon\operatorname{NE}(X)\hookrightarrow\operatorname{NE}(Y),
 \qquad\iota(X_k)=Y_{u+1+k}\quad(k\in[|X|]_{\Z}).
\]
Let
\begin{equation}\label{eq:shell-recolor}
 \gamma=\operatorname{Recol}(X,Y,\iota,g;f,\eta).
\end{equation}
Then, the map
\[
 \Phi_{\mathrm{sh}}\colon
 \prod_{s\in[j]_{\Z}}\mathcal C_{c_s}(X_s)
 \longrightarrow\mathcal C_{\sum_{s=1}^j c_s-\epsilon}(Y),
 \qquad \Phi_{\mathrm{sh}}\bigl((g_s,\gamma_s)_{s\in[j]_{\Z}}\bigr)=(g,\gamma)
\]
is injective. Also, we have
\[
 (g(\iota(X_k)),\gamma(\iota(X_k)))=(f(X_k),\eta(X_k))
 \qquad(k\in[|X|]_{\Z}).
\]
\item\label{c:shell-copy-precursor} For $i\in[j]_{\Z}$, $k\in[|X_i|]_{\Z}$, and $e=(X_i)_k$, define the copy of $e$ in $Y$ by
\begin{equation}\label{eq:shell-edge-copy}
 \widetilde e=Y_{r_{i-1}+k}.
\end{equation}
Then, we have
\begin{equation}\label{eq:shell-copy-precursor}
 P_{Y,g}(\widetilde w)=\widetilde{P_{X_i,g_i}(w)}
 \qquad(i\in[j]_{\Z},\ w\in N(X_i)).
\end{equation}
\item\label{c:shell-boundary-debts} For $i\in[0,j]_{\Z}$ and $e\in E(Y[0,r_i])$, we have
\[
 \operatorname{Debt}_{Y,g}(e,r_i)=
 \begin{cases}
 1,&e=Y_1,\\
 0,&e\ne Y_1.
 \end{cases}
\]
\item\label{c:shell-boundary-balances} We have
\[
 \operatorname{bBal}_{Y,g}(r_i)
 =1+i+\sum_{s=1}^i\operatorname{Surp}(X_s,g_s)
 \qquad(i\in[0,j]_{\Z}).
\]
\end{enumerate}
\end{lemma}
\begin{proof}
\prf{Claim~\ref{c:shell-bounded}} Let us verify that every shell grade lies in the prescribed range.
The bounds $u\in[m-1]_{\Z}$, $v\in[m-1]_{\Z}$, $j+1+\epsilon\le m$, and $g_i\in\operatorname{CtdG}_m(X_i)$ for $i\in[j]_{\Z}$ in
Definition~\ref{def:shell} give
\[
\begin{aligned}
 g_i(e)&\le m\qquad(i\in[j]_{\Z},\ e\in\operatorname{NE}(X_i)),\\
 \max\{u+1,v-1,j+1+\epsilon\}&\le m.
\end{aligned}
\]
Every edge of $Y$ inherits a grade from some $g_i$ or has grade
$0$, $u+1$, $v-1$, or $j+1+\epsilon$. Hence, $g\in G_m(Y)$.

\prf{Claim~\ref{c:shell-connected}} Let us verify the three conditions for connectedness in Definition~\ref{def:connected-grading}(1).
Let $y=\operatorname{Start}_{\mathrm{NE}}(Y)$ and $z=\operatorname{End}_{\mathrm{NE}}(Y)$.
We construct $P_{Y,g}$ in edge order, using Definition~\ref{def:stack}(1).
The grades $g(Y_1)=u+1$ and $g(Y_{k+1})=0$ for $k\in[u]_{\Z}$ give
\[
\begin{aligned}
 P_{Y,g}(Y_{k+1})&=y\qquad(k\in[u]_{\Z}),\\
 \operatorname{Debt}_{Y,g}(y,r_0)&=g(y)-u=(u+1)-u=1.
\end{aligned}
\]

Let us determine the h-precursors in each copy of $X_i$.
The subpath $Y[r_{i-1},r_i]$ is the translated copy of $X_i$ with the grading translated from $g_i$.
At each vertical edge $\widetilde w$ of this copy, the translated h-precursor from $X_i$ has positive debt and follows every horizontal edge before the copy.
Definition~\ref{def:stack}(1) therefore selects $\widetilde{P_{X_i,g_i}(w)}$ as the h-precursor of $\widetilde w$.

Let us compute the debt of $y$ after each copy of $X_i$.
Since $g_i$ is connected, Claim~\ref{c:node-contained} of Lemma~\ref{lem:node-structure}
shows that every horizontal edge of $X_i$ has debt zero at $|X_i|$.
Thus,
\[
 \operatorname{Debt}_{Y,g}(y,r_i)=1
 \qquad(i\in[j]_{\Z}).
\]

Let us compute the change in boundary balance across the copies.
The boundary balances satisfy
\[
 \operatorname{bBal}_{Y,g}(r_i)-\operatorname{bBal}_{Y,g}(r_{i-1})
 =|X_i|_E-\operatorname{TNG}_{X_i,g_i}
 =\operatorname{Surp}(X_i,g_i)+1
 \qquad(i\in[j]_{\Z}),
\]
and hence
\[
 \operatorname{bBal}_{Y,g}(r_j)-\operatorname{bBal}_{Y,g}(r_0)
 =\sum_{i=1}^j\operatorname{Surp}(X_i,g_i)+j.
\]

Let us prove $\operatorname{Node}_{Y,g}(y)=Y$.
The grades $g(Y_{r_j+1})=v-1$ and $g(Y_{r_j+1+k})=0$ for $k\in[v-1]_{\Z}$ give
\[
\begin{aligned}
 P_{Y,g}(Y_{r_j+1+k})&=Y_{r_j+1}\qquad(k\in[v-1]_{\Z}),\\
 \operatorname{bBal}_{Y,g}(r_j+v)-\operatorname{bBal}_{Y,g}(r_j)
 &=1-\sum_{k=1}^{v-1}g(Y_{r_j+1+k})=1,\\
 \operatorname{Debt}_{Y,g}(Y_{r_j+1},r_j+v)&=(v-1)-(v-1)=0.
\end{aligned}
\]
Immediately before $z$, the edge $y$ has debt one and every other horizontal edge has debt zero.
Hence,
\[
 P_{Y,g}(z)=y,\qquad
 \operatorname{Debt}_{Y,g}(y,|Y|)=0.
\]
This constructs $P_{Y,g}$ and gives $\operatorname{Node}_{Y,g}(y)=Y$.

Let us verify the node balance condition. For $i\in[j]_{\Z}$, $w\in N(X_i)$, and
$a=P_{X_i,g_i}(w)$, the identity $P_{Y,g}(\widetilde w)=\widetilde a$ and the connectedness of $g_i$ give
\[
\begin{aligned}
 &\operatorname{eBal}_{Y,g}(\widetilde w)-\operatorname{eBal}_{Y,g}(\widetilde a)\\
 &\qquad=\operatorname{eBal}_{X_i,g_i}(w)-\operatorname{eBal}_{X_i,g_i}(a)\ge0.
\end{aligned}
\]
For $Y_{k+1}$ with $k\in[u]_{\Z}$ and $Y_{r_j+1+k}$ with $k\in[v-1]_{\Z}$, we have
\[
\begin{aligned}
 \operatorname{eBal}_{Y,g}(Y_{k+1})&=1=\operatorname{eBal}_{Y,g}(y)
 &&(k\in[u]_{\Z}),\\
 \operatorname{eBal}_{Y,g}(Y_{r_j+1+k})&=\operatorname{eBal}_{Y,g}(Y_{r_j+1})
 &&(k\in[v-1]_{\Z}).
\end{aligned}
\]
By Lemma~\ref{lem:residual-nonnegative}, each $\operatorname{Surp}(X_i,g_i)$ is nonnegative.
Under condition~\ref{hyp:shell-preserving} or~\ref{hyp:shell-draining}, the edge balance at $z$ satisfies
\[
 \operatorname{eBal}_{Y,g}(z)=\operatorname{EGrDiff}_{Y,g}(y,z)
 =1+\sum_{i=1}^j\operatorname{Surp}(X_i,g_i)-\epsilon
 \ge1=\operatorname{eBal}_{Y,g}(y).
\]
Therefore, $g$ is connected. By Claim~\ref{c:shell-bounded},
$g\in\operatorname{CtdG}_m(Y)$.

\prf{Claim~\ref{c:shell-residual}}
By~\eqref{eq:residual},
\[
 \operatorname{Surp}(Y,g)=\operatorname{EGrDiff}_{Y,g}(y,z)-1
 =\sum_{i=1}^j\operatorname{Surp}(X_i,g_i)-\epsilon.
\]

\prf{Claim~\ref{c:shell-recolor}} Let us define the shell map and prove injectivity.
Claims~\ref{c:shell-connected} and~\ref{c:shell-residual} give
$(g,\gamma)\in\mathcal C_{\sum_{s=1}^j c_s-\epsilon}(Y)$.
Equation~\eqref{eq:shell} gives $g(\iota(X_k))=f(X_k)$.
Definition~\ref{def:recolor} then gives $\gamma(\iota(X_k))=\eta(X_k)$.
Restriction to each copy of $X_s$ recovers $(g_s,\gamma_s)$, proving injectivity.

\prf{Claim~\ref{c:shell-copy-precursor}} Let us compare the h-precursors in $X_i$ and its copy in $Y$.
Fix $i\in[j]_{\Z}$.
The copy map~\eqref{eq:shell-edge-copy} preserves edge order and grades.
For $w=(X_i)_k\in N(X_i)$, Definition~\ref{def:stack}(1) gives $\operatorname{Debt}_{X_i,g_i}(P_{X_i,g_i}(w),k-1)>0$.
Induction on $k$ preserves the debts of the horizontal edges in $X_i[0,k-1]$ under the copy map.
Every horizontal edge before the copy of $X_i$ precedes $\widetilde{P_{X_i,g_i}(w)}$.
Thus, Definition~\ref{def:stack}(1) selects $\widetilde{P_{X_i,g_i}(w)}$ as the h-precursor of $\widetilde w$.
Hence,~\eqref{eq:shell-copy-precursor} holds.

\prf{Claim~\ref{c:shell-boundary-debts}} Let us compute the debts at $r_i$.
The first $u$ vertical edges have h-precursor $Y_1$.
Hence,
\[
 \operatorname{Debt}_{Y,g}(Y_1,r_0)=g(Y_1)-u=1.
\]
By Claim~\ref{c:shell-copy-precursor}, every vertical edge in a copy of $X_s$ has its h-precursor in that copy.
Claim~\ref{c:node-contained} of Lemma~\ref{lem:node-structure} gives debt zero at $|X_s|$ for every horizontal edge of $X_s$.
Thus, at $r_i$, the edge $Y_1$ has debt $1$ and every horizontal edge in the copies of $X_1,\ldots,X_i$ has debt $0$.

\prf{Claim~\ref{c:shell-boundary-balances}} The definition of $\operatorname{bBal}$ and~\eqref{eq:residual} give
\[
 \operatorname{bBal}_{Y,g}(r_i)
 =1+\sum_{s=1}^i\bigl(|X_s|_E-\operatorname{TNG}_{X_s,g_s}\bigr)
 =1+i+\sum_{s=1}^i\operatorname{Surp}(X_s,g_s)
 \qquad(i\in[0,j]_{\Z}).
\]
\end{proof}

\begin{definition}[Fixed simple-run grading]\label{def:simple-run-grading}
Let $j\in[m]_{\Z}$.
Define $g_j^{\mathrm{run}}\in G_m(R_j)$ and $\gamma_j^{\mathrm{run}}\in\operatorname{Col}(R_j,g_j^{\mathrm{run}},m)$ by
\[
 g_j^{\mathrm{run}}(\operatorname{Start}_{\mathrm{NE}}(R_j))=j,\qquad
 g_j^{\mathrm{run}}(v)=0\quad(v\in N(R_j)),\qquad
 \gamma_j^{\mathrm{run}}(e)=1\quad(e\in\operatorname{NE}(R_j)).
\]
\end{definition}

\begin{lemma}\label{lem:word-bounds}
Let $U$ be a run with $\nu$ simple runs. Let $L$ be an integer.
Assume the following.
\begin{hypotheses}
\item\label{hyp:grading-range} $L\in[m-1]_{\Z}$.
\item\label{hyp:grading-lengths} Each simple run of $U$ has vertical length at most $L$.
\end{hypotheses}
Then, we have
\[
 m^{\ceil{\frac{\nu}{2}}}\le C_0(U)\le 2^{m(L+1)\nu}.
\]
\end{lemma}
\begin{proof}
Let us prove the lower bound first for one and two simple runs.
By Lemma~\ref{lem:word-reindexing}, we translate $U$ so that $\operatorname{Start}_{\mathrm{pt}}(U)=(0,0)$.

First, let us prove the bound for $\nu=1$.
Let $U=R_j$ and let $g=g_j^{\mathrm{run}}$.
This grading is connected and satisfies $\operatorname{Surp}(U,g)=0$.
By assumptions~\ref{hyp:grading-range} and~\ref{hyp:grading-lengths},
$1\le j\le L\le m-1$. Thus, the horizontal edge has $\binom mj\ge m$ colors.

Second, let us prove the bound for $\nu=2$.
Let $U=R_u\ccn R_v$ and use
$g=\operatorname{GShell}_{\emptyset}(;u,v,0)$. Claims~\ref{c:shell-connected} and~\ref{c:shell-residual} of Lemma~\ref{lem:shell} give a connected grading of surplus zero.
The edge $\operatorname{End}_{\mathrm{NE}}(U)$ has grade one and $m$ colors.

Let us extend the lower bound by induction on $\nu$.
For $\nu\ge3$, let $U=R_u\ccn X\ccn R_v$ with $|X|_E=\nu-2$.
For each $(g,\gamma)\in\mathcal C_0(X)$, Claims~\ref{c:shell-connected}--\ref{c:shell-recolor} of Lemma~\ref{lem:shell}, followed by Lemma~\ref{lem:word-reindexing}, give a colored grading $(h,\delta)\in\mathcal C_0(U)$ such that
\begin{equation}\label{eq:word-bound-restriction}
 h(U_{u+1+k})=g(X_k),\qquad
 \delta(U_{u+1+k})=\gamma(X_k)\quad(k\in[|X|]_{\Z}).
\end{equation}
The edge $\operatorname{End}_{\mathrm{NE}}(U)$ has grade two and $\binom m2\ge m$ colors.
Equation~\eqref{eq:word-bound-restriction} recovers $(g,\gamma)$ from $(h,\delta)$.
Induction gives
\[
 C_0(U)\ge\binom m2 C_0(X)
 \ge m\,m^{\ceil{\frac{\nu-2}{2}}}=m^{\ceil{\frac{\nu}{2}}}.
\]

Let us prove the upper bound by summing the choices of grade and color on each edge.
By assumption~\ref{hyp:grading-lengths}, there are at most $(L+1)\nu$ edges. Hence,
\[
 C_0(U)\le\prod_{e\in\operatorname{NE}(U)}\sum_{r=0}^m\binom mr
 =2^{m|U|}\le2^{m(L+1)\nu}.
\]
\end{proof}

\begin{lemma}\label{lem:outer-shell-extension}
Let $X$ be a tNE path. Let $A$ and $B$ be runs with $P\ge1$ and $S\ge1$
simple runs, respectively.
Assume the following.
\begin{hypotheses}
\item\label{hyp:outer-shell-lengths} Each simple run of $A$ and $B$ has vertical length at most $m-1$.
\item\label{hyp:outer-shell-counts} $\max(P,S)\le2\min(P,S)$.
\end{hypotheses}
Then, for every integer $d\ge0$, there is an injective map
\[
 i_d\colon\mathcal C_d(X)\longrightarrow\mathcal C_d(A\ccn X\ccn B)
\]
with the following property. For each $(g,\gamma)\in\mathcal C_d(X)$ and $(h,\delta)=i_d(g,\gamma)$, we have
\begin{equation}\label{eq:outer-extension-restriction}
 h\bigl((A\ccn X\ccn B)_{|A|+k}\bigr)=g(X_k),\qquad
 \delta\bigl((A\ccn X\ccn B)_{|A|+k}\bigr)=\gamma(X_k)
 \quad(k\in[|X|]_{\Z}).
\end{equation}
\end{lemma}
\begin{proof}
Let us construct shell gradings that preserve the surplus $d$.
Let $(g,\gamma)\in\mathcal C_d(X)$. Let $u\in[m-1]_{\Z}$, $v\in[m-1]_{\Z}$, and $t\in[m-1]_{\Z}$.
Give $R_t$ the colored grading $(g_t^{\mathrm{run}},\gamma_t^{\mathrm{run}})$.
The shell paths
\begin{equation}\label{eq:outer-shells}
\begin{gathered}
 \operatorname{Shell}(X;u,v,0),\\
 \operatorname{Shell}(R_t,X;u,v,0),\\
 \operatorname{Shell}(X,R_t;u,v,0)
\end{gathered}
\end{equation}
add $(1,1)$, $(2,1)$, and $(1,2)$ simple runs before and after $X$, respectively.
For each shell path $Y$, let $h$ be its shell grading from~\eqref{eq:shell}.
Claims~\ref{c:shell-connected} and~\ref{c:shell-residual} of Lemma~\ref{lem:shell} give $h\in\operatorname{CtdG}_m(Y)$ and
\[
 \operatorname{Surp}(Y,h)=\operatorname{Surp}(X,g)=d.
\]

Let us choose the numbers of shell paths to add exactly $P$ simple runs before $X$ and $S$ after $X$.
By assumption~\ref{hyp:outer-shell-counts}, the following decomposition has nonnegative integer coefficients:
\begin{equation}\label{eq:outer-shell-counts}
 (P,S)=
 \begin{cases}
 (P-S)(2,1)+(2S-P)(1,1),& S\le P\le2S,\\
 (S-P)(1,2)+(2P-S)(1,1),& P\le S\le2P.
 \end{cases}
\end{equation}

Let us match the simple runs of $A$ and $B$ in order.
Let $a_1,\ldots,a_P$ and $b_1,\ldots,b_S$ be the ordered vertical lengths
of the simple runs of $A$ and $B$.
By assumption~\ref{hyp:outer-shell-lengths},
\[
 a_i\in[m-1]_{\Z}\quad(i\in[P]_{\Z}),\qquad
 b_i\in[m-1]_{\Z}\quad(i\in[S]_{\Z}).
\]

Let us specify one step of the construction.
Let $p\in[0,P]_{\Z}$ and $q\in[0,S]_{\Z}$.
Suppose that $Y$ is a translated copy of
$R_{a_{p+1}}\ccn \cdots\ccn R_{a_P}\ccn X\ccn R_{b_1}\ccn \cdots\ccn R_{b_q}$.
For the three shell operators in~\eqref{eq:outer-shells} applied to $Y$, use
\[
 (u,v)=(a_p,b_{q+1}),\qquad
 (u,t,v)=(a_{p-1},a_p,b_{q+1}),\qquad
 (u,t,v)=(a_p,b_{q+1},b_{q+2}),
\]
respectively.
The first choice requires $p\ge1$ and $q<S$.
The second requires $p\ge2$ and $q<S$.
The third requires $p\ge1$ and $q+2\le S$.

Let us complete the construction.
Construct the binary shell paths before the unary shell paths.
Starting at $(p,q)=(P,0)$, the numbers of shell paths
in~\eqref{eq:outer-shell-counts} give $(p,q)=(0,S)$.
At each step, define the shell grading by~\eqref{eq:shell} and its coloring by~\eqref{eq:shell-recolor}.
At $(p,q)=(0,S)$, $Y$ is a translated copy of $A\ccn X\ccn B$.
Lemma~\ref{lem:word-reindexing} transports the colored grading on $Y$ to $A\ccn X\ccn B$.

Let us verify~\eqref{eq:outer-extension-restriction} and prove that $i_d$ is injective.
By Claim~\ref{c:shell-recolor} of Lemma~\ref{lem:shell}, the map $i_d$ copies the grade and color of $X_k$ to $(A\ccn X\ccn B)_{|A|+k}$ for every $k\in[|X|]_{\Z}$.
This proves~\eqref{eq:outer-extension-restriction}.
For $(g_1,\gamma_1)\in\mathcal C_d(X)$ and $(g_2,\gamma_2)\in\mathcal C_d(X)$,
equation~\eqref{eq:outer-extension-restriction} gives
\[
 i_d(g_1,\gamma_1)=i_d(g_2,\gamma_2)
 \ \Longrightarrow\ 
 \begin{cases}
 g_1(X_k)=g_2(X_k),\\
 \gamma_1(X_k)=\gamma_2(X_k)
 \end{cases}
 \quad(k\in[|X|]_{\Z}).
\]
Hence, $(g_1,\gamma_1)=(g_2,\gamma_2)$ and $i_d$ is injective.
\end{proof}

\begin{lemma}[A shifted supermultiplicative limit]\label{lem:shifted-growth}
Let $(A_n)_{n\in\Z_{>0}}$ be a sequence of real numbers and let $C\ge0$.
Assume
\[
 A_n\ge1,\qquad \log A_n\le Cn\quad(n\in\Z_{>0}),\qquad
 A_{i+j+2}\ge A_iA_j\quad(i,j\in\Z_{>0}).
\]
Then, the following statements hold.
\begin{enumerate}
\item\label{c:shifted-growth-limit} The limit exists and satisfies
\[
 L:=\lim_{n\to\infty}\frac{\log A_n}{n}
   =\sup_{n\ge1}\frac{\log A_n}{n+2}\in[0,C].
\]
\item\label{c:shifted-growth-bound} For $n\in\Z_{>0}$, we have $A_n\le e^{(n+2)L}$.
\end{enumerate}
\end{lemma}
\begin{proof}
\prf{Claim~\ref{c:shifted-growth-limit}} Let us write $\nu=n+2$.
Let $B_\nu=\log A_{\nu-2}$ for $\nu\in\Z_{\ge3}$ and let
\[
 L_0=\sup_{\nu\ge3}\frac{B_\nu}{\nu}.
\]
We have $0\le B_\nu\le C\nu$ for $\nu\in\Z_{\ge3}$ and
$B_{u+v}\ge B_u+B_v$ for $u\in\Z_{\ge3}$ and $v\in\Z_{\ge3}$.
Thus, $L_0\in[0,C]$.
Fix $t\in\Z_{\ge3}$.
For $\nu\ge3$, let $\nu=kt+r$ for $k\in\Z_{\ge0}$ and $r\in[3,t+2]_{\Z}$.
Induction on $k$ gives
\[
 B_\nu\ge kB_t+B_r\ge kB_t.
\]
Since $\frac{k}{\nu}\to\frac1t$, we obtain
\[
 \liminf_{\nu\to\infty}\frac{B_\nu}{\nu}\ge\frac{B_t}{t}.
\]
Taking the supremum over $t$ gives a lower bound of $L_0$.
The definition of $L_0$ gives $\frac{B_\nu}{\nu}\le L_0$.
Hence, $\frac{B_\nu}{\nu}\to L_0$ and
\[
 \frac{\log A_n}{n}=\frac{n+2}{n}\frac{B_{n+2}}{n+2}\longrightarrow L_0.
\]
Replacing $\nu$ by $n+2$ in the definition of $L_0$ proves Claim~\ref{c:shifted-growth-limit}.

\prf{Claim~\ref{c:shifted-growth-bound}} Let us bound the error in the finite growth estimate.
Claim~\ref{c:shifted-growth-limit} gives $\log A_n\le(n+2)L$.
Exponentiation gives the assertion.
\end{proof}

\begin{proposition}[Periodic growth]\label{prop:periodic}
Let $W$ be a run with $q$ simple runs. Assume the following.
\begin{hypotheses}
\item\label{hyp:periodic-lengths} Every simple run of $W$ has vertical length at most $m-1$.
\item Let $A_j=C_0(W^j)$ for $j\ge1$.
\end{hypotheses}
Then, we have the following.

\begin{enumerate}
\item\label{c:periodic-supermultiplicativity} For all integers $i\ge1$ and $j\ge1$, we have
\[
 A_{i+j+2}\ge A_iA_j.
\]
\item\label{c:periodic-limit} We have
\begin{equation}\label{eq:periodic-limit}
 \lim_{j\to\infty}\frac{\log A_j}{j}
       =\sup_{j\ge1}\frac{\log A_j}{j+2}
       \in(0,\infty).
\end{equation}
\item\label{c:periodic-bound} For every $j\ge1$, we have
\[
 A_j\le\exp\left((j+2)\lim_{k\to\infty}\frac{\log A_k}{k}\right).
\]
\end{enumerate}
\end{proposition}
\begin{proof}
\prf{Claim~\ref{c:periodic-supermultiplicativity}} Let us construct a colored grading on $W^{i+j+2}$ from colored gradings on $W^i$ and $W^j$.
Let $i\ge1$ and $j\ge1$.
Let $a_1,\ldots,a_q$ be the vertical lengths of the simple runs of $W$.
Let $X=W^i$ and $Y=W^j$.
For $(g_1,\gamma_1)\in\mathcal C_0(X)$ and $(g_2,\gamma_2)\in\mathcal C_0(Y)$, form the shell paths
\[
\begin{aligned}
 Z_1&=\operatorname{Shell}(X,Y;a_q,a_1,0),\\
 Z_{s+1}&=\operatorname{Shell}(Z_s;a_{q-s},a_{s+1},0)
 \quad(s\in[q-1]_{\Z}).
\end{aligned}
\]
Give each $Z_s$ its shell grading~\eqref{eq:shell} and coloring~\eqref{eq:shell-recolor}.
Assumption~\ref{hyp:periodic-lengths} and Claims~\ref{c:shell-connected} and~\ref{c:shell-residual} of Lemma~\ref{lem:shell} give a connected colored grading of surplus zero on $Z_q$.
By Lemma~\ref{lem:word-reindexing}, translate $Z_q$ and its colored grading to
$Z=W\ccn X\ccn Y\ccn W=W^{i+j+2}$. Denote the resulting colored grading by $(h,\delta)$.
This defines the map
\[
 \Phi_{i,j}\colon\mathcal C_0(W^i)\times\mathcal C_0(W^j)
 \longrightarrow\mathcal C_0(W^{i+j+2}),\qquad
 \Phi_{i,j}((g_1,\gamma_1),(g_2,\gamma_2))=(h,\delta).
\]

Let us recover $(g_1,\gamma_1)$ and $(g_2,\gamma_2)$ to prove that $\Phi_{i,j}$ is injective.
Claim~\ref{c:shell-recolor} of Lemma~\ref{lem:shell} gives
\begin{equation}\label{eq:periodic-restrictions}
\begin{aligned}
 h(Z_{|W|+k})&=g_1(X_k),&
 \delta(Z_{|W|+k})&=\gamma_1(X_k)
 &&(k\in[|X|]_{\Z}),\\
 h(Z_{|W|+|X|+k})&=g_2(Y_k),&
 \delta(Z_{|W|+|X|+k})&=\gamma_2(Y_k)
 &&(k\in[|Y|]_{\Z}).
\end{aligned}
\end{equation}
Equation~\eqref{eq:periodic-restrictions} determines both $(g_1,\gamma_1)$ and $(g_2,\gamma_2)$ from $(h,\delta)$.
Hence, $\Phi_{i,j}$ is injective and $A_{i+j+2}\ge A_iA_j$.

\prf{Claim~\ref{c:periodic-limit}} Let us prove~\eqref{eq:periodic-limit}.
Lemma~\ref{lem:word-bounds} gives $A_j\ge1$ and $\log A_j\le m^2qj\log2$ for $j\ge1$.
Claim~\ref{c:periodic-supermultiplicativity} gives $A_{i+j+2}\ge A_iA_j$ for $i\in\Z_{>0}$ and $j\in\Z_{>0}$.
Claim~\ref{c:shifted-growth-limit} of Lemma~\ref{lem:shifted-growth} therefore gives~\eqref{eq:periodic-limit}.
The lower bound in Lemma~\ref{lem:word-bounds} also gives
\[
 \lim_{j\to\infty}\frac{\log A_j}{j}\ge\frac q2\log m>0.
\]

\prf{Claim~\ref{c:periodic-bound}}
Claim~\ref{c:shifted-growth-bound} of Lemma~\ref{lem:shifted-growth} proves Claim~\ref{c:periodic-bound}.
\end{proof}

\begin{definition}\label{def:periodic-growth-rate}
Let $W$ be a run whose simple runs have vertical lengths at most $m-1$.
We define
\[
 \ell(W)=\lim_{j\to\infty}\frac{\log C_0(W^j)}{j}.
\]
\end{definition}

\subsection{A uniform upper bound for finite dimension pairs}
\begin{lemma}[Reduction of dimension pairs]\label{lem:dimension-reduction}
Let $m\in\Z_{\ge3}$ and let $a\in\Z_{>0}$ and $b\in\Z_{>0}$ be coprime.
Let $Q=mab-a^2-b^2>0$.
Then, there exist coprime $u\in\Z_{>0}$ and $v\in\Z_{>0}$ such that
\[
\begin{gathered}
 u\le v\le\frac{mu}{2},\qquad muv-u^2-v^2=Q,\qquad E_m(u,v)=E_m(a,b),\\
 u+v\le\sqrt{\frac{m+2}{m-2}}\sqrt Q.
\end{gathered}
\]
\end{lemma}
\begin{proof}
Let us obtain $(u,v)$ by decreasing $x+y$.
Let
\[
 (x,y)=(\min(a,b),\max(a,b)).
\]
If $2y>mx$, replace $(x,y)$ by
\[
 (\min(x,mx-y),\max(x,mx-y)).
\]
Claims~\ref{c:moduli-symmetries-identities}, \ref{c:moduli-coprimality}, and~\ref{c:moduli-quadratic}
of Lemma~\ref{lem:moduli-symmetries} show that each replacement preserves
the Euler characteristic, coprimality, and $Q$.
The inequality $mxy-x^2-y^2=Q>0$ implies $y<mx$.
Thus, both new entries are positive.
Their sum changes by $mx-2y<0$.
A strictly decreasing sequence of positive integers is finite. Thus, the replacements end at a pair
$u\le v\le\frac{mu}{2}$.

Let us bound $u+v$.
We have $mv-2u\ge0$ and $mu-2v\ge0$.
Therefore,
\[
 (m+2)Q-(m-2)(u+v)^2=(mv-2u)(mu-2v)\ge0.
\]
Dividing by $m-2>0$ and taking square roots gives $u+v\le\sqrt{\frac{m+2}{m-2}}\sqrt Q$.
\end{proof}

\begin{proposition}\label{prop:finite-endpoint}
For every $(a,b)\in\cA_m$, let $\rho(a,b)=\frac{b}{a}$.
Then, we have
\begin{equation}\label{eq:finite-endpoint}
\begin{split}
 0\le\frac{\log E_m(a,b)}a
 &\le m\log2\sqrt{\frac{m+2}{m-2}}
             \sqrt{\max(m\rho(a,b)-1-\rho(a,b)^2,0)}\\
 &\hspace{12mm}+\frac{\log\left(\frac{2^m}{m}\right)}a.
\end{split}
\end{equation}
\end{proposition}
\begin{proof}
Let us separate the cases according to the sign of $Q=mab-a^2-b^2$.
First, suppose $Q>0$.
Lemma~\ref{lem:dimension-reduction} gives a coprime pair $(u,v)$ satisfying
\[
 u\le v\le\frac{mu}{2},\qquad E_m(u,v)=E_m(a,b),\qquad
 u+v\le\sqrt{\frac{m+2}{m-2}}\sqrt Q.
\]

Let us bound the Euler characteristic at the reduced pair $(u,v)$.
If $u\ge2$, we have $v>u$ because $u$ and $v$ are coprime. Then, Claim~\ref{c:frame-dictionary} of Proposition~\ref{prop:frame} applies.
Equation~\eqref{eq:ambient} gives $\Gamma(u,v)<\frac vu\le\frac m2$.
The floor-difference formula in Definition~\ref{def:frame} gives
\[
 \Lambda_i(u,v)\le\ceil{\Gamma(u,v)}\le\ceil{\frac m2}\le m-1
 \qquad(i\in[u]_{\Z}).
\]
The final simple run is $R_1$ by~\eqref{eq:frame-word}.
The lower bound in Lemma~\ref{lem:word-bounds}
with $L=m-1$ and $\nu=u+1$ gives $C_0(\mathcal U(u,v))\ge m$.

Let us obtain the upper bound when $u\ge2$.
The run $\mathcal U(u,v)$ has exactly $u+v+1$ edges.
At each edge, the number of choices of a grade and its color is
\[
 \sum_{k=0}^m\binom{m}{k}=2^m.
\]
Therefore,
\[
 C_0(\mathcal U(u,v))\le(2^m)^{u+v+1}=2^{m(u+v+1)}.
\]
Dividing by $m$ in~\eqref{eq:dictionary} gives
\begin{equation}\label{eq:finite-reduced}
 1\le E_m(u,v)\le\frac{2^{m(u+v+1)}}m.
\end{equation}
If $u=1$, the stability condition~\eqref{eq:stability} means that the images of the $m$ arrow maps span $\C^v$.
The moduli space $\cM^s_{(1,v)}(K(m))$ is the Grassmannian of $v$-dimensional quotients of
$\C^m$. Then, $E_m(1,v)=\binom mv$. The bounds in~\eqref{eq:finite-reduced} also hold for $u=1$.

Let us apply $u+v\le\sqrt{\frac{m+2}{m-2}}\sqrt Q$ from Lemma~\ref{lem:dimension-reduction}.
Taking logarithms in~\eqref{eq:finite-reduced} therefore yields
\begin{equation}\label{eq:finite-Q}
 0\le\log E_m(a,b)
 \le m\log2\sqrt{\frac{m+2}{m-2}Q}+\log\left(\frac{2^m}{m}\right).
\end{equation}

Second, suppose $Q=0$. No positive integral pair has this value, because $m\ge3$ gives
\[
 (m-1)^2<m^2-4<m^2.
\]

Third, suppose $Q<0$. Let us prove $Q=-1$ and $E_m(a,b)=1$.
The nonempty stable moduli space $\cM^s_{(a,b)}(K(m))$ has dimension
$1+Q$~\cite[part (2) of Proposition 4.3]{Weist}. Since $Q$ is a negative integer, we obtain $Q=-1$.

Let us prove $E_m(a,b)=1$ when $Q=-1$.
The stable locus is a nonempty open subset of $\operatorname{Hom}_{\C}(\C^a,\C^b)^m$ and is therefore irreducible.
The smooth geometric quotient $\cM^s_{(a,b)}(K(m))$ is irreducible and zero-dimensional. Hence, $\cM^s_{(a,b)}(K(m))$ consists of one point.
Replacing $Q$ by $\max(Q,0)$ in~\eqref{eq:finite-Q} gives a bound on all of $\cA_m$.
Dividing by $a$ gives~\eqref{eq:finite-endpoint}.
\end{proof}


